% ==============================================================================
% Section 13 & References: Econometric Limitations & Academic Bibliography
% Target: Page 15
% ==============================================================================

\section{Econometric Limitations and Concluding Recommendations}
\label{sec:limitations_conclusions}

\subsection{Econometric Limitations of the Study}

While the empirical analyses conducted in this report provide rigorous quantitative benchmarks for managing equity market risk, several econometric limitations must be acknowledged:
\begin{enumerate}[leftmargin=1.5em, itemsep=1pt, parsep=0pt, label=(\alph*)]
    \item \textbf{Assumption of Stationarity:} Non-parametric Historical Simulation implicitly assumes that the underlying data-generating process is strictly stationary. Structural breaks—such as the corporate demerger of ITC Hotels or major revisions in the GST compensation cess structure—alter the true distribution of future asset returns.
    \item \textbf{Ghost Features and Sample Horizon Sensitivity:} In unweighted historical windows, a large historical shock (e.g., an $8.5\%$ Budget drawdown) abruptly enters the sample, inflates $\VaR$ estimates for exactly 246 days, and then abruptly drops out, creating an artificial ``ghost drop'' in estimated risk without any corresponding change in market fundamentals.
    \item \textbf{Unhedged Cash Equity Horizon:} The models analyzed assume static positions over 1-day and 10-day intervals, ignoring dynamic intraday stop-loss execution, exchange circuit filters, and derivatives-based tail risk hedging.
\end{enumerate}

\subsection{Concluding Recommendations for Institutional Desks}

Based on our empirical findings across the 246-session observation window of ITC Limited:
\begin{itemize}[leftmargin=1.5em, itemsep=1pt, parsep=0pt]
    \item \textbf{Adopt Expected Shortfall over VaR:} Given that 10-day $99\%$ Expected Shortfall (\inr 95,131) exceeds 10-day $99\%$ Historical $\VaR$ (\inr 61,488) by over $54.7\%$, institutions must anchor internal capital reserves to $\ES$ to guarantee solvency against deep tail events.
    \item \textbf{Incorporate Policy Jump-Diffusion Processes:} Standard Gaussian Monte Carlo understates extreme tail risk. Institutional simulation desks should augment Geometric Brownian Motion with Poisson jump processes calibrated specifically to Union Budget event dates.
    \item \textbf{Execute Derivative Tail Hedging:} Holding an unhedged single-stock exposure requires a punitive regulatory capital reserve of \inr 1,84,464 ($37.5\%$ of position value). Implementing option-based hedging structures offers substantial capital relief and tail insulation.
\end{itemize}

\vspace{-0.4em}

% ==============================================================================
% Peer-Reviewed Academic Bibliography
% ==============================================================================
\begingroup
\footnotesize
\setlength{\bibsep}{1.2pt}
\renewcommand{\section}[2]{}% Suppress default bibliography section header
\subsubsection*{References}
\bibliographystyle{plainnat}
\bibliography{references}
\endgroup
